\chapter{Reproducibility and Supplementary Diagnostics}
\section{Immutable Evidence}
The report is based on the code/results release at revision:
\begin{quote}\small\path{50f8995a014f4d68714cba7a4ea81e911c859374}\end{quote}
The original training revision is:
\begin{quote}\small\path{2d0dfbf34480bba5e602f165fd12dc0c3ef75ef8}\end{quote}
Later report commits do not
rewrite the historical training record.

The main result CSV contains 3,840 rows and has SHA-256
\begin{quote}\small\ttfamily
\path{2576804e3caadd3025ff95a9835729ced219fd38dce4ca30ee4a399727a61150}.
\end{quote}
The archive \path{releases/code-results-50f8995.tar.gz} has SHA-256
\begin{quote}\small\ttfamily
\path{73b15ce8a707f671f7d7102d8ac776d504a9786afda5d933f0a72e7363861e58}.
\end{quote}
It contains 264 source and artifact files, including the 181 items in the
frozen audit inventory. The source repository is
\url{https://github.com/Ovais09/csi6900-noisy-labels}. Git LFS stores the large
archive separately from the ordinary Git pointer. It must be fetched before
extracting the release. The archive must not be unpacked over an active
working tree, and untrusted serialized models must not be loaded.

\section{Artifact Map}
\begin{longtable}{>{\raggedright\arraybackslash}p{.38\linewidth}>{\raggedright\arraybackslash}p{.55\linewidth}}
\toprule
Artifact & Purpose \\
\midrule
\endhead
\path{configs/experiment_config.yaml} & Exact experimental grid. \\
\path{results/tables/main_experiment_results.csv} & Per-combination and per-condition metrics and counts. \\
\path{results/tables/summary_results.csv} & 960 rows aggregating the four seeds within settings. \\
\path{results/tables/analysis/} & Eight analysis tables, including paired cleaning and direct detector comparisons. \\
\path{results/models/} & Four RF models for one Dry Bean setting, not all models fitted by the benchmark. \\
\path{results/figures/} & 60 main and 13 diagnostic figures, each in PNG and PDF. \\
\path{results/logs/run_manifest.json} & Config, data, environment, result, and original training provenance. \\
\path{results/logs/release_manifest.json} & Artifact sizes and SHA-256 hashes. \\
\path{results/logs/final_audit.json} & Passing audit of that frozen inventory. \\
\path{report/generated/evidence.json} & Checksummed inputs used to generate report tables and figures. \\
\bottomrule
\end{longtable}

The report generator recomputes analysis tables from the exact CSV and checks
them against the frozen versions before producing LaTeX tables. Numeric
summary tables are not manually transcribed. Report-specific plots average
the same matrix; selected diagnostic PDFs are copied unchanged after hash
verification. Their source values and revision are recorded in the evidence
file. The IEEEtran BibTeX style produces the numbered bibliography.

\section{Build and Reproduction Commands}
The complete extraction and audit commands are in \path{releases/README.md}.
The following commands assume the Git-LFS archive has been verified and
restored to an empty \path{releases/restored/} directory:
\begin{lstlisting}
python scripts/make_report_assets.py --release-root releases/restored
docker build -f report/Dockerfile -t csi6900-report .
docker run --rm --network none -v "$PWD:/work" csi6900-report
\end{lstlisting}
The last command runs pdfLaTeX and BibTeX through latexmk and writes
\path{report/build/main.pdf}. Building from the already committed report
assets does not require model training or downloading datasets. A full
benchmark rerun uses the archived configuration and pipeline in its own
checkout; it should not overwrite the frozen evidence used by this report.

\section{Preparation and Review}
AI assistance was used during software implementation, testing, analysis
support, and manuscript preparation. The numerical evidence is tied to the
saved experiment artifacts, not generated as illustrative data. Source
verification and software checks do not replace the author's responsibility
to understand the methods, review the conclusions, and comply with the
supervisor's and institution's requirements for assistance and disclosure.

\section{Supplementary Macro F1 and Curves}
Figure~\ref{fig:f1} provides the macro-F1 counterpart of the cyclic accuracy
plot. Figures~\ref{fig:roc} and \ref{fig:pr} concern only the saved noisy RF
classifier from the Dry Bean diagnostic setting. Their smooth ranking
summaries should not be interpreted as proof that cleaning improves the
hard-label predictions in that setting.

\begin{figure}[htbp]
\centering\includegraphics[width=\linewidth]{f1_macro_class_dependent.pdf}
\caption{Macro F1 under cyclic class-dependent corruption. Each point averages
three models and four seeds within a dataset. The plotted methods and rates
match the main accuracy comparison.}
\label{fig:f1}
\end{figure}
\begin{figure}[htbp]
\centering\includegraphics[width=.9\linewidth]{noisy_baseline_roc_curve.pdf}
\caption{Micro-averaged ROC curve for the noisy Dry Bean RF classifier,
symmetric 20\% corruption and seed 42. Micro ROC AUC is approximately 0.9878;
the main metric for this model is macro one-vs-rest AUROC, approximately 0.9874.}
\label{fig:roc}
\end{figure}
\begin{figure}[htbp]
\centering\includegraphics[width=.9\linewidth]{noisy_baseline_precision_recall_curve.pdf}
\caption{Micro-averaged precision--recall curve for the same noisy classifier.
Trapezoidal PR AUC is approximately 0.9558. This is not average precision.}
\label{fig:pr}
\end{figure}
